\section{Horns on rooted trees}\label{sec:tree-horns}
Let $\Omega$ be the dendroidal category. Its objects are rooted trees
and its morphisms are morphisms of the free symmetric coloured operads
on those trees \cite[\S3]{MoerdijkWeissDendroidal}. For a tree $T$ with
edge set $E(T)$, put $I(T)=\prod_{a=1}^k(E(T)\sqcup C_a)$, where each
$C_a$ is a fixed finite set. Pullback along colour maps, extended by
the identity on the $C_a$, defines $D:\Omega^{\mathrm{op}}\to
\mathbf{Mod}_A$, with $D(T)=A^{I(T)}$. On linear trees the colour maps
are the ordinal maps, so this restricts to $X(\mathbf c;A)$ for
$c_a=|C_a|$.

An elementary inner face contracts an inner edge. For a tree with
at least two vertices, an elementary outer face prunes an endpoint
of the tree of vertices, retaining its one incident inner edge.
The faces of a one-vertex tree are its individual edges. For a face
$F:S\to T$, the horn omitting $F$ is the union of the representables
of the other elementary faces of $T$. Its data are tensors on those
faces agreeing on common subfaces. This description follows by
factoring dendroidal morphisms into degeneracies, isomorphisms and
faces \cite[Lemma~2.3]{BasicNikolaus}.
Write $K_F(T)$ for the tensors that vanish on the included faces and
$K(T)$ for those that vanish on all elementary faces. Identify face
index sets with their images in $I(T)$.

\begin{proposition}\label{prop:tree-lifts}
Let $F:S\to T$ be an elementary face. The pullback
$F^*:K_F(T)\to K(S)$ has kernel $K(T)$ and image
\[
 A^{\,I(S)\setminus\bigcup_{G\ne F} I(S_G)}
 \ \subseteq\ K(S),\qquad G:S_G\to T.
\]
The inclusion is by extension by zero in $I(S)$. The map onto this
image splits by extension by zero in $I(T)$.
\end{proposition}
\begin{proof}
Each elementary face of $S$, composed with $F$, factors through
another elementary face of $T$ by the codimension-two face identity
\cite[Lemma~2.1]{BasicNikolaus}. Thus the restriction of $K_F(T)$
lies in $K(S)$. Membership in $K_F(T)$ means vanishing on the union
of the included index sets. Restriction to $I(S)$ and extension by zero
give the stated image and section. The kernel vanishes on $F$
as well as on the included faces.
\end{proof}

Take the two-vertex tree with vertices $u:a\to e$ and
$v:(e,b)\to r$, and omit the contraction of $e$. Its two included faces have colour
sets $\{a,e\}$ and $\{e,b,r\}$, meeting in the single-edge tree on $e$.
Every such inner horn fills by gluing the two arrays on their common
index set. With two axes and empty tails, $S$ is the corolla
$(a,b)\to r$ and $K(S)$ is spanned by its six off-diagonal units.
The omitted-face map reaches the four ordered pairs joining $a$
to $b$ or $r$. The units at $(b,r)$ and $(r,b)$ are visible on the
face with colours $\{e,b,r\}$ and have no lifts in $K_F(T)$.

\begin{theorem}\label{thm:tree-cutoff}
For $k$ axes and arbitrary fixed tails, all horns on trees with at
least $k+2$ vertices fill uniquely. For each $k\geq2$, there is an
inner horn on a tree with $k+1$ vertices whose restriction has
cokernel free of rank $k!$. For $k=1$, every inner horn fills uniquely.
\end{theorem}
\begin{proof}
For an ambient index tuple, join two included faces that contain it
when a common subface contains it. Compatible data assign one value
to each connected component of this graph, whereas a filler assigns
one value to all copies of the tuple. Hence a horn fills for all
compatible data if every nonempty graph is connected.

Let $T$ have $n\geq3$ vertices and $o$ outer faces. The outer vertices
are the endpoints of the tree of vertices, so $o\geq2$. The included
contractions form a clique among the faces containing a tuple, as
do the included prunings. A contraction and a pruning also give
adjacent copies when the contracted edge is away from the pruned
vertex. In each case the two moves commute, and the common face
retains their common colours. If a tuple has a disconnected graph,
exactly two included faces retain it. They are the pruning of an
outer vertex $v$ and the contraction of its incident inner edge $e$.
Indeed, an outer vertex has only one incident inner edge, and for
$n\geq3$ an inner edge has at most one outer endpoint.

The remaining included faces delete pairwise disjoint sets of
edges. The tuple must contain at least one edge from each such set,
so it contains at least $n+o-4$ distinct edges. If $o=2$, the tree
of vertices is a path. Pruning $v$ and then its neighbour, or
contracting $e$ and then pruning the merged vertex, gives the same
common face. It contains the edges deleted by every other included
contraction and, if the other pruning is included, all the edges
deleted by that pruning. A disconnected tuple must contain an
additional edge outside this common face. In either case it contains
at least $n-1$ distinct edges. Thus $n\geq k+2$ gives existence.
For uniqueness, a coordinate of $K_F(T)$ must hit the deletion
sets of all $n+o-2$ included faces, so $K_F(T)=0$ if $n>k$.

For sharpness, take $k+1$ vertices on a path, all unary except the
second vertex from the top, which has an extra leaf $b$. Let $e$
be the inner edge incident to the top vertex, and omit contraction
of another inner edge $h$. Use a tuple consisting of the root
edge, $b$, and every inner edge except $e$ and $h$. Only contraction
of $e$ and pruning of the top vertex retain this tuple. On
contracting $e$ an operation using the merged vertex has the
top leaf as an input, and that leaf is absent from the pruned tree.
On pruning the top vertex an operation using the second vertex
has $e$ as an input, and $e$ is absent from the contracted tree.
The colour $b$ is therefore isolated from the root in their common
subfaces. The two copies of the tuple may take distinct values,
which prevents a filler. For this tree the obstructing tuples are
exactly the permutations of the one just described. At the top
endpoint the forced inner edges and the root require one further
colour, which must be $b$. At the root endpoint, if its incident
contraction is included, the common face contains every shared
colour, since the neighbour of the root vertex is unary. The
cokernel of the horn restriction is therefore free of rank $k!$.

For one axis the bound covers all trees with at least three
vertices. On a two-vertex tree the two faces in the inner horn
meet in their connecting edge, so the horn fills by gluing.
Their deletion sets give uniqueness.
\end{proof}
